\chapter{Objetivos}\label{cap:objetivos}

El objetivo del presente Trabajo Fin de Grado consiste en desarrollar y evaluar diferentes modelos de Inteligenia Artificial mediante aprendizaje supervisado, capaces de predecir valores de referencia (set points) que optimicen el funcionamiento de redes eléctricas de distribución. El propósito de estos modelos será predecir de manera efectiva y precisa estos valores de referencia, reduciendo la necesidad de disponer de información detallada sobre la topología y parámetros eléctricos de la red, requeridos por métodos tradicionales de optimización como el algoritmo Optimal Power Flow. Con este fin, el entrenamiento de estos modelos se realizarán a partir de grandes volúmenes de datos medidos por sensores instalados en la red.

Para alcanzar este objetivo general, se plantean los siguientes objetivos específicos:
\begin{enumerate}
	\item \textbf{Análisis e Investigación de Tecnologías:} Seleccionar las librerías y frameworks de Inteligencia Artificial a utilizar y realizar un análisis e investigación exhaustiva de estas tecnologías, para una correcta implementación y optimización de los modelos predictivos en el contexto eléctrico.
	\item \textbf{Ingeniería de Datos y Preprocesamiento:} Tratar el dataset disponible mediante tareas de limpieza, normalización y transformación, aplicando técnicas de selección de características (Feature Selection), para reducir la dimensionalidad y garantizar la calidad de los datos de entrada de los modelos.
	\item \textbf{Construcción y Entrenamiento de Modelos Predictivos:} Desarrollar y entrenar arquitecturas de aprendizaje supervisado basadas en Random Forest, XGBoost y Redes Neuronales que se adapten a la problemática.
	\item \textbf{Optimización de Hiperparámetros:} Aplicar estrategias de optimización de hiperparámetros sobre los modelos implementados, para maximizar métricas de rendimiento y evitar el sobreajuste (overfitting).
	\item \textbf{Evaluación y Análisis Comparativo:} Evaluar el desempeño de los modelos utilizando métricas estadísticas rigurosas sobre conjuntos de datos de prueba no vistos anteriormente por los modelos, para realizar una comparación técnica entre estos y permitir una reflexión sobre su aplicabilidad real en contextos eléctricos.
\end{enumerate}